\bookchapter{第二十七章：父母亦能輕鬆用——適老化手機完全指南}

\chapterintro{曰句實話，寫此一章之時，吾欲至了我爸媽。我媽每次手機彈出一個“系統更新”，她皆不敢點，怕將手機弄壞矣。我爸之微信字體比螞蟻尚小，但他不知道何以調大，遂此麼眯着眼睛觀了兩年。此非彼等“笨”，乃手機默認以所有人皆當成矣二十多歲之年輕人。此一章，吾用最直白則，以長輩等手機使用中最頭疼之問題一個個解開。}

\booksection{零、先以心態擺正：此非“老人用不來”，乃“手機沒作好”}

於作任何設置之前，有一個觀念要先曰清楚：\textbf{而今之智能手機界面，自設計之初遂沒何以考慮老年用戶}。彈窗小得觀不清、操作手勢須精準觸碰、圖標長得皆差不多——此非用戶之問題，乃產品之問題。故帶長輩用手機，第一件事非教，乃幫彼等\textbf{以環境調好}。環境調好了，大半之問題遂消失矣。

帶長者以手機有三條鐵律：(1)\textbf{勿催}——汝用矣十年之智能機，彼等或方用矣三個月。(2)\textbf{用類較}——“Wi-Fi遂像家裏之水管，開矣水龍頭方有水”，此種比喻比專業術語好理解一萬倍。(3)\textbf{寫下來}——口頭教一遍彼等或忘矣，汝拿張紙將關鍵步驟寫下來貼於他手機殼背面，比何皆管用。

\booksection{一、首步：開“簡易模式”——手機變清爽矣}

而今之手機基本皆內置矣一個“簡易模式”（有之叫“關懷模式”“老年模式”“長輩模式”），打開之後，整個手機皆將大變樣：

\begin{itemize}
  \item \textbf{字體與圖標變得甚大}——再亦不用戴老花鏡觀手機矣。
  \item \textbf{桌面圖標大幅精簡}——唯保留電話、短信、相機、相冊、微信等最常用之幾個App，其他全收起來。
  \item \textbf{操作變簡單矣}——甚多花裏胡哨之功能（負一屏、智能助手、浮窗等）默認關閉，界面主打一個清爽。
  \item \textbf{聲音變大矣}——來電鈴聲、通知音量、通話聲音皆自動調至最大。
\end{itemize}

\booksubsection{各品牌簡易模式開啓方法（直接照着操作遂行）：}

\begin{center}
\begin{tabularx}{\linewidth}{|X|X|X|}
\hline
\textbf{品牌} & \textbf{模式名稱} & \textbf{開法} \tabularnewline
\hline
iPhone & 輔助用（無單獨簡易模式，但可逐項調大） & 設置 > 顯示與亮度 > 文字大小，拉至最大。再進"輔助功能 > 顯示與文字大小"開"粗體文本""按鈕形狀" \tabularnewline
電荷泵 & 簡易模式 & 設置 > 系統與更新 > 簡易模式 > 開啓 \tabularnewline
小米/紅米 & 至簡模式 & 設 > 益多設 > 至簡模式 > 開啓 \tabularnewline
OPPO/一加 & 簡易模式 & 設 > 特色能 > 簡易模式 \tabularnewline
vivo/iQOO & 簡易桌面 & 設 > 桌面、鎖屏與壁紙 > 簡易桌面 \tabularnewline
三星 & 簡易模式 & 設 > 顯 > 簡易模式 \tabularnewline
\hline
\end{tabularx}
\end{center}

若汝用之乃其他品牌，於手機“設置”頂部之搜索框裏直接搜“簡易”或“長輩”兩個字，一般皆能尋到。尋不至則，按下面之“手動調大”方案一條條設置，效果亦一樣。

\booksection{二、若不開簡易模式，手調亦可——“大字+大圖標”五步走}

有些長輩不欲用簡易模式，覺得界面太“老氣”。那遂手動調以下五項，效果一樣好：

\begin{enumerate}
  \item \textbf{字體調至最大}：設置 > 顯示 > 字體大小，拖至最右邊。順便以"顯示大小"（有之叫"顯示縮放"）亦調至最大，此樣所有界面元素皆與着變大。
  \item \textbf{圖標變大}：於主屏幕雙指捏合（像縮小照片那樣）進入桌面設置 > 桌面佈局（或圖標大小），選格子最少之佈局（比如3×5），圖標自然遂大矣。
  \item \textbf{開"高對較度文字"}：使文字邊緣更銳利，於"輔助功能"裏可尋到此開關。
  \item \textbf{關閉不必要之通知}：設置 > 通知，以除矣電話、短信、微信之外之App通知全關掉。此樣手機未能一將兒響一下，未能使人覺得煩。
  \item \textbf{以常用聯繫人放至桌面}：長按桌面空白處 > 添加小工具 > 聯繫人，以老伴、兒女、醫生、緊急聯繫人之快捷撥號放至桌面上。點一下遂撥出去，不用翻通訊錄。
\end{enumerate}

\begin{tcbpractice}{動手：助長者"裝修"一下手機}
\begin{enumerate}
  \item 打開簡易模式（按上面之品牌表格操作）。
  \item 以電量數字顯示打開（設置 > 電池 > 電量百分比），使彼等隨時知道尚剩幾何電。
  \item 調出通知欄，以藍牙、NFC、飛行模式此長輩用不着之快捷開關移除或藏起來。
  \item 於桌面上放四個快捷撥號：老伴、兒女、家庭醫生（若有）、社區/物業。
\end{enumerate}

\end{tcbpractice}

\booksection{三、防詐——此章請必帶着長者共觀}

老年人乃電信詐騙之重災區。非彼等警惕性低，而乃騙子則術經過矣精心設計，專門針對老年人之心理弱點。以下乃吾整理之最常見騙局類型與應對方法，建議汝直接念賜長輩聽，或打印出來貼於電話旁邊。

\booksubsection{六大常見局（聞之矣即刻掛電話）：}

\begin{itemize}
  \item \textbf{(1) "汝孫子出事矣，趕緊匯錢"}——騙子將曰"汝孫子/孫女於外面闖禍矣，爲警察抓矣，須保證金"。記住：\textbf{任何要求轉賬才能"放人"之電話皆乃詐騙}。公檢法未能電話要錢。
  \item \textbf{(2) "汝之醫保卡爲停用矣"}——短信裏曰醫保卡異常，使汝點鏈接"激活"。記住：\textbf{醫保局、社保局未能發短信使汝點鏈接}。遇至此種情況直接打12333（社保官方熱線）覈實。
  \item \textbf{(3) "高益理財，穩賺不賠"}——微信羣、短信裏薦各種"養舊理財""國家扶持項目"，承諾高額報。切記：\textbf{承諾年化收益超過5\%尚"保本"之，99.9\%乃騙局}。
  \item \textbf{(4) "吾乃汝領導/老朋友，換號碼矣"}——冒充熟人借。切記：\textbf{打電話覈實比何皆管用}，曰"不方便接電話"之八成乃騙子。
  \item \textbf{(5) "免費領雞蛋/保健品講座"}——線下免費活動誘導老人聽"健康講座"，而後賣天價保健品。記住：\textbf{天上不掉餡餅，免費活動背後皆乃生意}。
  \item \textbf{(6) "汝之手機中毒矣，須遠程修復"}——騙子使汝下載一個遠程控制App（如向日葵、TeamViewer），而後遠程操控汝之手機轉走錢。記住：\textbf{絕對勿使任何人遠程操控汝之手機！}手機廠商、通信運營商、銀行皆未能賜汝打電話使汝裝遠程控制軟件。
\end{itemize}

\booksubsection{三道防線，立好之後能擋掉90\%之詐騙：}

\begin{enumerate}
  \item \textbf{首道：手機自帶攔截}——安卓手機打開"騷擾攔截"（於撥號App或手機管家App裏），iPhone打開"靜音未知來電"（設置 > 電話 > 靜音未知來電）。
  \item \textbf{次道：運營商攔截}——發送短信"KTFSR"至10086，免費開通移動之高頻騷擾電話攔截；電信聯通也有類似服務。打運營商客服電話使人工幫汝開通遂行。
  \item \textbf{三道：國家反詐中心App}——於應用商店下載"國家反詐中心"App，註冊並開啓來電預警。此App與公安之數據庫聯動，騙子號碼一撥過來遂能識別。
\end{enumerate}

\begin{tcbwarn}{賜長輩之"防騙口訣"（打印出來貼電話旁邊）}
\begin{itemize}
  \item \textbf{陌生電話要轉賬 掛掉！}
  \item \textbf{短信鏈接要汝點 刪掉！}
  \item \textbf{曰高收益保本之 拉黑！}
  \item \textbf{使汝裝遠程控制軟件 直接掛！}
  \item \textbf{真警察未能電話辦案更未能要錢！}
  \item \textbf{拿不準之事，先賜兒女打個電話問問。}
\end{itemize}

\end{tcbwarn}

\booksection{四、遠程協助——汝身於異地亦能助父母修手機}

有時候長輩遇至之問題，汝於電話裏描述半天彼等亦尋不至對應之按鈕。此時候“遠程協助”功能遂派上用場矣——汝可直接於自己之手機上看到彼等手機之屏幕，幫彼等操作。各品牌之方案不一樣：

\begin{itemize}
  \item \textbf{電荷泵}：自帶"暢連"App，裏面有"屏幕共享"功能。汝與爸媽皆裝暢連，發起視頻通話後點擊"共享屏幕"，汝就可於彼等屏幕上畫圈標註，告訴汝媽媽"點此裏"。
  \item \textbf{小米/紅米}：小米視頻通話App裏有"遠程協助"，支屏幕共享+語音通話同步。
  \item \textbf{OPPO/vivo}：內置"家庭空間"或"遠程守護"，於設置裏搜索"遠程"遂能尋到。
  \item \textbf{通用方案（所有品牌皆可用）}：下載"向日葵遠程控制"App（免費版夠用），汝等各裝一個。註冊同一個賬號或加好友，使爸媽將手機上顯示之驗證碼念賜汝，汝遂能遠程看到與操作彼等之屏幕矣。
  \item \textbf{iPhone之用戶}：FaceTime通話中支持"共享屏幕"。於FaceTime通話時，點擊屏幕共享按鈕，對方遂能看到汝之屏幕。但iOS安全限制無法遠程操作，只能"指揮式教學"。
\end{itemize}

\booksection{五、微信之“適舊改”——與長者曰此幾句遂夠矣}

微信乃長輩等用得最多之App，但界面太密太雜。以下乃幾個立竿見影之調整：

\begin{itemize}
  \item \textbf{調大微信自己之字體}：微信 > 吾 > 設置 > 通用 > 字體大小，拖至最大。注意，此唯影響微信內部之字體，不影響手機系統字體。
  \item \textbf{開"關懷模式"}：微信 > 吾 > 設置 > 關懷模式 > 開啓。打開後字體與按鈕皆將變大，色彩對比度加強，還能一鍵"聽文字消息"（點一下文字消息遂能朗讀出來）。
  \item \textbf{幫彼等備註好重要聯繫人}：以兒女、親戚之名字於通訊錄裏改成容易認之——比如"大兒子-小明""二女兒-小紅"，使彼等一打開微信遂知道誰乃誰。
  \item \textbf{教彼等用"語音輸入"而非打字}：聊天框旁邊之鍵盤上，按住空格鍵（安卓）或點麥克風圖標（iPhone），直接說話遂能變成文字發送。識別準確率已甚高矣，對未能拼音打字之老人來曰乃"質變級"之便利。
  \item \textbf{關閉"發現頁"裏不用之功能}：微信 > 吾 > 設置 > 通用 > 發現頁管理，以"朋友圈""視頻號""直播""購物""遊戲"裏不用之全關掉，減少幹擾。
\end{itemize}

\booksection{六、手機丟矣何以辦？助長者作個“應急預案”}

預則立，不預則廢。趁手機尚沒丟，先以此幾件事作好：

\begin{enumerate}
  \item \textbf{確認"查尋吾之手機"已打開}：iPhone於設置 > 汝之名字 > 查尋 > 查尋吾之iPhone；安卓於設置搜"查尋設備"。務必確認此開關乃亮着之。
  \item \textbf{以重要之密碼與賬號記於一張紙上}：Apple ID / 華爲賬號 / 小米賬號之密碼，鎖於家裏抽屜裏一份，賜信得過之兒女一份。丟矣手機之後要登錄此賬號來遠程鎖定與定位，若忘矣密碼遂麻煩矣。
  \item \textbf{教長者一句話}：\textbf{"手機丟矣，第一時間賜家人打電話。"}家人可立刻登錄iCloud或安卓之查尋設備網頁，遠程標記爲"丟失之模"，此樣撿至之人亦打不開、用不了。
\end{enumerate}

\booksection{七、無線耳機與助聽器兼容——爲忽視之好能}

甚多長輩有聽力問題，但不知道手機可與助聽器協同工作。iPhone與部分安卓手機支持MFi（Made for iPhone）助聽器，可直接以電話聲音傳至助聽器裏，比舉着手機聽筒清楚太多矣。

\begin{itemize}
  \item \textbf{連法}：iPhone設置 > 輔助功能 > 助聽設備，打開藍牙，按照助聽器之配對說明操作。
  \item \textbf{若不用助聽器}：於"輔助功能 > 音頻/視覺"裏打開"電話降噪"與"單聲道之音"，能改善通話清晰度。
  \item \textbf{安卓}：於"輔助功能 > 聽力增強"裏有類似設置，各品牌位置略有差異，直接搜"助聽"兩個字。
\end{itemize}

\booksection{八、健康相關能——手機不僅乃通訊工具}

\begin{itemize}
  \item \textbf{SOS緊急求助}：連續按電源鍵五次（大多數手機），手機會自動撥打緊急電話並賜預設之緊急聯繫人發送定位。幫長輩提前設置好——iPhone："設置 > 緊急SOS"，安卓："設置 > 安全 > 緊急求助"。教長輩用之前先演習一遍，免得到時候緊張按不準。
  \item \textbf{跌倒檢測}：Apple Watch有此功能，檢測至汝摔倒且一段時間不動，將自動撥打急救電話並通知家人。有預算則賜長輩配一個，關鍵時刻能救命。
  \item \textbf{用藥提醒}：iPhone之"健康"App內置用藥提醒功能，設好藥名、時間與劑量，到點遂彈出提醒。安卓可用"定時鬧鐘"或下載"喫藥提醒"類App。
\end{itemize}

\begin{tcbsummary}{本章之要}
\begin{itemize}
  \item 先開簡易模式或手動調大字體/圖標，環境調好大半問題遂沒矣
  \item 帶長輩用手機三鐵律：勿催、用類比、寫下來
  \item 六大常見局要屢提醒，三道防線立好能擋90\%詐
  \item 遠程協助使汝人於外地也能幫爸媽修手機
  \item 微信開"關懷模式"+"語音輸入"乃長輩最實用之兩個功能
  \item SOS緊急求助與跌倒檢測，關鍵時刻救命之配置值得花五分鐘設好
\end{itemize}

\end{tcbsummary}
